\centerline{\bf \Huge 7 класс}

\problem{\textbf{1}} \textbf{Критерии:} \\
\textit{7 баллов} $-$ Найден любой верный пример.\\
\textit{0 баллов} $-$ Любое неверное решение.\\
\textit{0 баллов} $-$ Решение отсутствует.\\
\textit{0 баллов} $-$ При обоснованных подозрениях на списывание участником из любого источника (Интернет, учебное пособие, ключи олимпиады и т.д.) член жюри имеет право указать на такие фрагменты с помощью восклицательных знаков на полях. \\


\problem{\textbf{2}} \textbf{Критерии:} \\
\textit{0-4 балла} $-$ Решение А.\\
\textit{0-1 балл} $-$ Составление текстового условия Б.\\
\textit{0-2 балла} $-$ Любое правильное решение Б.\\
\textit{0 баллов} $-$ Любой ответ без объяснений.\\
\textit{0 баллов} $-$ Решение отсутствует.\\
\textit{0 баллов} $-$ При обоснованных подозрениях на списывание участником из любого источника (Интернет, учебное пособие, ключи олимпиады и т.д.) член жюри имеет право указать на такие фрагменты с помощью восклицательных знаков на полях. \\


\problem{\textbf{3}} \textbf{Критерии:} \\
\textit{7 баллов} $-$ Любое полное верное решение.\\
\textit{6-7 баллов} $-$ Верное решение. Имеются небольшие недочёты, в целом не 
влияющие на решение.\\
\textit{5-6 баллов} $-$ Решение в целом верное. Однако оно содержит ошибки, либо пропущенные случаи, не влияющие на логику рассуждений.\\
\textit{0-4 баллов} $-$ Имеются продвижения с объяснениями.\\
\textit{0 баллов} $-$ Любой ответ без объяснений.\\
\textit{0 баллов} $-$ Решение отсутствует.\\
\textit{0 баллов} $-$ При обоснованных подозрениях на списывание участником из любого источника (Интернет, учебное пособие, ключи олимпиады и т.д.) член жюри имеет право указать на такие фрагменты с помощью восклицательных знаков на полях. \\

\problem{\textbf{4}} \textbf{Критерии:} \\
\textit{7 баллов} $-$ Любое полное верное решение.\\
\textit{5-6 баллов} $-$ Любое верное решение, но доказательство того, что могут быть только два ответа неполное.\\
\textit{4 балла} $-$ Найдены два правильных ответа с обоснованием, без попыток доказательства невозможности других.\\
\textit{2 балла} $-$ Найден один правильных ответ с обоснованием, без попыток доказательства невозможности других.\\
\textit{0 баллов} $-$ Любой ответ без объяснений.\\
\textit{0 баллов} $-$ Решение отсутствует.\\
\textit{0 баллов} $-$ При обоснованных подозрениях на списывание участником из любого источника (Интернет, учебное пособие, ключи олимпиады и т.д.) член жюри имеет право указать на такие фрагменты с помощью восклицательных знаков на полях. \\

\problem{\textbf{5}} \textbf{Критерии:} \\
\textit{7 баллов} $-$ Любое полное верное решение.\\
\textit{2 балла} $-$ Новый квадрат должен быть $5\times5$, так как $4\cdot4+3\cdot3=25$\\
\textit{0 баллов} $-$ Любой ответ без объяснений.\\
\textit{0 баллов} $-$ Решение отсутствует.\\
\textit{0 баллов} $-$ При обоснованных подозрениях на списывание участником из любого источника (Интернет, учебное пособие, ключи олимпиады и т.д.) член жюри имеет право указать на такие фрагменты с помощью восклицательных знаков на полях. \\
